\begin{table}[htbp]
\centering
\caption{The six core non-MHC Tier 1 genes with colocalization and fine-mapping support.}
\label{tab:tier1}
\small
\resizebox{\textwidth}{!}{%
\begin{tabular}{lllllllll}
\toprule
  Gene & Chr & PPH4 & $n_{\text{CS}}$ & PIP & $p_{\text{SMR}}$ & Known drug & Risk categories & High-risk phenotypes \\
\midrule
  UTP11  & 1 & 0.9996 & 10 & 1.00 & $1.56\times10^{-13}$ & --- & 4 & 6 \\
  BMP2   & 20 & 0.9971 & 10 & 1.00 & $5.77\times10^{-15}$ & --- & 2 & 2 \\
  PNKD   & 2 & 0.9337 & 2 & 1.00 & $2.38\times10^{-15}$ & --- & 4 & 8 \\
  LAMC1  & 1 & 0.9294 & 10 & 1.00 & $7.14\times10^{-24}$ & OCRIPLASMIN (approval) & 3 & 3 \\
  TMBIM1 & 2 & 0.8262 & 3 & 1.00 & $5.03\times10^{-14}$ & --- & 3 & 7 \\
  GPBAR1 & 2 & 0.8132 & 3 & 1.00 & $1.75\times10^{-14}$ & --- & 2 & 8 \\ 
\bottomrule
\end{tabular}
}

\tblnote{Colocalization is defined as PPH4 $>$ 0.8 and fine-mapping support as convergence of SuSiE to a credible set. Risk categories and high-risk phenotypes are from Phenome-Wide Association Study (PheWAS) screening in the Open Targets Platform (see Methods and Supplementary Table~S8): risk categories = number of distinct organ-system categories carrying at least one genetic association for the gene; high-risk phenotypes = number of non-CRC phenotypes with genetic evidence in predefined high-risk categories.}
\end{table}
